\documentclass[11pt]{article}
\usepackage{calc,fancyhdr,lastpage}
\usepackage{listings}
\usepackage[hmargin=.75in,vmargin=1in,
            footskip=.55in,headsep=.55in-\headheight]{geometry}
\usepackage{amsmath,amssymb}
\usepackage{hyperref}
\usepackage{graphicx}


% Formats, symbols, abbreviations.
\let\altemph\textsl
\let\strong\textbf
\let\code\texttt
\let\latinabb\emph
\newcommand*{\etc}{\latinabb{etc}}
\newcommand*{\eg}{\latinabb{e.g.}}
\newcommand*{\ie}{\latinabb{i.e.}}
% To get proper-looking symbols in \texttt.
\newcommand*{\txtbksl} {\symbol{"5C}}% \
\newcommand*{\txtcaret}{\symbol{"5E}}% ^
\newcommand*{\txtunder}{\symbol{"5F}}% _
\newcommand*{\txtlcurl}{\symbol{"7B}}% {
\newcommand*{\txtrcurl}{\symbol{"7D}}% }
\newcommand*{\txttilde}{\symbol{"7E}}% ~

% Commands \question[marks]{title} and \subquestion[marks]{title}.
\newcounter{questionnumber}
\newcommand*{\question}[2][]
   {\refstepcounter{questionnumber}\section*
    {Problem \thequestionnumber.\quad
        \ifx\empty#1\empty\else[#1 marks]\quad\fi#2}}
\newcounter{subquestionnumber}[questionnumber]
\renewcommand*{\thesubquestionnumber}{\alph{subquestionnumber}}
\newcommand*{\subquestion}[2][]
   {\refstepcounter{subquestionnumber}\subsubsection*
    {Part (\thesubquestionnumber)\quad
       \ifx\empty#1\empty\else[#1 marks]\quad\fi#2}}

% Redefine `enumerate' to use less vertical space.
\let\etaremune\enumerate
\let\etaremunedne\endenumerate
\renewenvironment{enumerate}
   {\etaremune
    \setlength{\topsep}{.25ex plus .125ex minus .1825ex}%
    \setlength{\itemsep}{\topsep}\setlength{\parsep}{0ex}%
    \setlength{\leftmargin}{1.75em}\setlength{\labelsep}{.5em}%
    \setlength{\labelwidth}{1.75em}\ignorespaces}
   {\etaremunedne}

% Redefine `itemize' to use less vertical space.
\let\ezimeti\itemize
\let\ezimetidne\enditemize
\renewenvironment{itemize}
   {\ezimeti
    \setlength{\topsep}{.25ex plus .125ex minus .1825ex}%
    \setlength{\itemsep}{\topsep}\setlength{\parsep}{0ex}%
    \setlength{\leftmargin}{1.75em}\setlength{\labelsep}{.5em}%
    \setlength{\labelwidth}{1.75em}\ignorespaces}
   {\ezimetidne}

%% A heading in the instructions.
\newcommand*{\heading}[1]{\subsubsection*{#1}}

% Headings.
\pagestyle{fancy}
\let\headrule\empty
\let\footrule\empty
\lhead{{\bfseries ANONYMOUS COURSE}}
\chead{{\bfseries\large Lab}}
\rhead{{\bfseries Fall 2024}}
\lfoot{{Anonymous Department}}
\cfoot{{}}
\rfoot{{Page \thepage\ of \pageref{LastPage}}}


\begin{document}
\lstset{language=Python}

\section{HOPFIELD NETWORKS}

The following code is handed out:

\begin{lstlisting}[language=Python]
def E(x0, x1, x2, w01, w02, w12):
    term1 = x0 * x1 * w01
    term2 = x0 * x2 * w02
    term3 = x1 * x2 * w12
    return -(term1 + term2 + term3)

def print_all_energies(w01, w02, w12):
    for x0 in [-1, 1]:
        for x1 in [-1, 1]:
            for x2 in [-1, 1]:
                print("x: (", x0, x1, x2, ")", "E:", E(x0, x1, x2, w01, w02, w12))
\end{lstlisting}

This could be used with the weights from last week as follows:

\begin{lstlisting}[language=Python]
w01 = 2
w02 = -1
w12 = 1
print_all_energies(w01, w02, w12)
\end{lstlisting}

The output is as follows:

\begin{lstlisting}
x: ( -1 -1 -1 ) E: -2
x: ( -1 -1 1 ) E: -2
x: ( -1 1 -1 ) E: 4
x: ( -1 1 1 ) E: 0
x: ( 1 -1 -1 ) E: 0
x: ( 1 -1 1 ) E: 4
x: ( 1 1 -1 ) E: -2
x: ( 1 1 1 ) E: -2
\end{lstlisting}

Recall that the minima of the energy function are ``memories''. This means that the network ``remembers'' the values $(-1, -1, -1)$, $(-1, -1, 1)$, $(1, 1, -1)$, $(1, 1, 1)$. We would now like to ``store'' a new memory in the network. We will do it by adjusting the weights $w$ a little bit.

The strategy for storing a new memory $(x_0, x_1, x_2)$ is as follows:

\begin{enumerate}
    \item If $x_0 \cdot x_1 > 0$, increasing $w_{01}$ will decrease the energy of the network at $(x_0, x_1, x_2)$, so we will increase $w_{01}$ by $0.1$. Otherwise, we will decrease $w_{01}$ by $0.1$.
    \item If $x_0 \cdot x_2 > 0$, increasing $w_{02}$ will decrease the energy of the network $(x_0, x_1, x_2)$, so we will increase $w_{02}$ by $0.1$. Otherwise, we will decrease $w_{02}$ by $0.1$.
    \item If $x_1 \cdot x_2 > 0$, increasing $w_{12}$ will decrease the energy of the network at $(x_0, x_1, x_2)$, so we will increase $w_{12}$ by $0.1$. Otherwise, we will decrease $w_{12}$ by $0.1$.
\end{enumerate}

\subsection*{Part (a)}
Explain why the claim ``If $x_0 \cdot x_1 > 0$, increasing $w_{01}$ will decrease the energy of the network at $(x_0, x_1, x_2)$'' makes sense.

\subsection*{Part (b)}
Write three if-statements that execute the strategy above, and verify that the effect of that is to decrease the energy of the network at $(-1, 1, 1)$ when $x = (-1, 1, 1)$. Note: all we are saying is ``write three if-statements and use \texttt{print\_all\_energies()} to verify that the energy at $(-1, 1, 1)$ is decreased.

\subsection*{Part (c)}
Write a for-loop that repeatedly adjust the weights as above, and uses \texttt{print\_all\_energies()} Suggestion: \texttt{print "=============="} in between iterations to make the output more readable.

\subsection*{Part (d)}
Use the for-loop to verify that after about 4 iteration, the new memory is stored.

\subsection*{Part (e)}
Put the for-loop in a function. The function should take in the weights and the memory to store, and should return the weights (as a list with three elements) after the memory is stored.

\subsection*{Part (f)}
Use the function to store the memory $(1, -1, 1)$ in the network. Note that the function returns a list. There are several ways to extract the answers from a list. Use tools from the course (i.e., 0-th element of \texttt{L} is \texttt{L[0]}) to update $w_{01}$, $w_{02}$, $w_{12}$.


\end{document}
